\documentclass[10pt]{article} 

\usepackage[a4paper,margin=2cm,left=2cm,right=2cm]{geometry} 
\usepackage[UKenglish]{babel}

\usepackage{amsmath}
\usepackage{amssymb}
\usepackage{amsthm}

\newtheorem{theorem}{Theorem}
\newtheorem{lemma}{Lemma}
\newtheorem{corollary}{Corollary}[theorem]
\newtheorem{proposition}{Proposition}
\theoremstyle{definition}
\newtheorem{definition}{Definition}
\newtheorem{example}{Example}
\theoremstyle{remark}
\newtheorem{remark}{Remark}
\newtheorem{notation}{Notation}

\newcommand{\R}{\mathbb{R}}
\newcommand{\Z}{\mathbb{Z}}
\newcommand{\N}{\mathbb{N}}
\newcommand{\Q}{\mathbb{Q}}
\newcommand{\C}{\mathbb{C}}
\newcommand{\F}{\mathbb{F}}

% \DeclareMathOperator{\dom}{dom}

% \usepackage{graphicx}
% \graphicspath{{./images/}}

\title{Linear Algebra}
\author{joel}
\date{}

\begin{document}

\maketitle

\section*{Intro}

The notes here are based off of the book "Linear Algebra Done Right". 

\section{Vector Spaces}

\begin{notation}
    $\F$ stands for either $\R$ or $\C$. That is, if a theorem holds for $\F$, it holds
    for both the real and complex numbers. 
\end{notation}

The elements of $\F$ are called `scalars'. Lists of length $n$ are ordered
collections of $n$ objects. Cannot be infinite, must have the same order to be
classed as equal. Also often called an $n$-tuple. 

\begin{definition}
    $\F^n$ is the set of all lists of length $n$ of elements of $\F$: 
    \[ \F^n = \{ (x_1, \ldots, x_n) : x_k \in \F \text{ for } k = 1, \ldots, n\}. \]
    For $(x_1, \ldots, x_n) \in \F$ and $k \in \{1, \ldots, k\}$ we say that
    $x_k$ is the $k$th coordinate of $(x_1, \ldots, x_n)$.
\end{definition}

Addition is commutative in $\F^n$. 
There exist additive inverses in $\F^n$.
Scalar multiplication in $\F^n$ is such that we have $\lambda \in \F$ and $x \in \F^n$
to get $\lambda x \in F^n$.

\begin{definition}
    A vector space is a set $V$ along with an addition on $V$ and a scalar
    multiplication on $V$ such that we have commutativity, associativity, additive identity, 
    additive inverse, multiplicative identity, and distributive properties. 
\end{definition}


\begin{example}
    An example of a vector space involves a set of functions: 
    \begin{itemize}
        \item If $S$ is a set, then $\F^S$ denotes the set of functions from $S$ to $\F$. 
        \item For $f, g \in \F^S$, the sum $f + g \in \F^S$ is the function defined by 
            \[ (f +g)(x) = f(x) + g(x), \] for all $x \in S$. 
        \item For $\lambda \in \F$ and $f \in \F^S$, the product $\lambda f \in \F^S$ is the function defined by 
            \[ (\lambda f)(x) = \lambda f ( x), \] for all $x \in S$. 
    \end{itemize}
    For instance, if $\F = \R$ and $S = [0, 1]$ then $R^{[0,1]}$ is the set of
    all real valued functions on the interval $[0, 1]$.
\end{example}

A vector space has a unique additive identity. 
Every element in a vector space has a unique additive inverse. 

\end{document}
